% META: {"id": "TEXP-CAG-CO-008", "chapitre": "TEXP-COMPLEXES-ALGEBRE-GEOMETRIE", "type_objet": "corrige", "exercice_ref": "TEXP-CAG-EX-008", "capacites_codes": ["C3"], "sources_inspiration": [], "status": "generated"}
% BEGIN-VERIFY
% from sympy import I, symbols, conjugate, expand, simplify
% a, b, c, d = symbols('a b c d', real=True)
% z, w = a + I * b, c + I * d
% assert simplify(conjugate(z * w) - conjugate(z) * conjugate(w)) == 0
% for n in range(1, 7):
%     assert simplify(conjugate(z ** n) - conjugate(z) ** n) == 0
% assert simplify(conjugate(1 / z) - 1 / conjugate(z)) == 0
% assert expand(z * w) == expand((a * c - b * d) + I * (a * d + b * c))
% END-VERIFY
\begin{corrige}{TEXP-CAG-EX-008}
\textbf{1.} $zz'=(aa'-bb')+\mathrm{i}(ab'+a'b)$ donc
$\overline{zz'}=(aa'-bb')-\mathrm{i}(ab'+a'b)$. Par ailleurs
$\overline{z}\,\overline{z'}=(a-\mathrm{i}b)(a'-\mathrm{i}b')
=(aa'-bb')-\mathrm{i}(ab'+a'b)$. Les deux coïncident.

\textbf{2.} Au rang $1$ c'est immédiat. Si $\overline{z^n}=\overline{z}^{\,n}$,
alors $\overline{z^{n+1}}=\overline{z^n\cdot z}
=\overline{z^n}\cdot\overline{z}=\overline{z}^{\,n}\overline{z}
=\overline{z}^{\,n+1}$ d'après la question 1.

\textbf{3.} De $z\times\frac1z=1$ on tire, en conjuguant et en utilisant la
question 1, $\overline{z}\times\overline{\left(\frac1z\right)}=\overline{1}=1$,
donc $\overline{\left(\frac1z\right)}=\frac{1}{\overline{z}}$.
\end{corrige}
